\section{Depart from Me}

At the end of the year the Church keeps the feast of Christ the King. In Year A of the Postconciliar Missal its Gospel is the scene that closes our Lord's last discourse in Matthew: the Son of man come in his majesty with all the angels, the nations gathered before him, the sheep set on his right hand and the goats on his left.\endnote{Matthew 25:31--46, the Gospel of the Solemnity of our Lord Jesus Christ, King of the Universe, in Year A of the Postconciliar Missal. In the 1962 Missal the same passage is read on the Monday of the first week of Lent.} To those on his right he says, ``Come, ye blessed of my Father, possess you the kingdom prepared for you from the foundation of the world.'' Then he turns to the others:

\begin{quote}
Then he shall say to them also that shall be on his left hand: Depart from me, you cursed, into everlasting fire, which was prepared for the devil and his angels.\endnote{Matthew 25:34 and 25:41, Douay--Rheims (Challoner).}
\end{quote}

The sentence is spoken by the one who said, ``Come to me all you that labor and are burdened, and I will refresh you,''\endnote{Matthew 11:28, Douay--Rheims (Challoner).} and both sayings are his. The Lord himself chose the words of the sentence, the Church has kept them as he gave them, and no preacher may soften them.

Saint John Chrysostom, preaching to his people at Antioch, noticed how the sentence is framed. The blessed are sent to a kingdom ``prepared for you''; the cursed are sent to a fire prepared for someone else.

\begin{quote}
Even here He shows His love, and indicates the yearning which He hath toward thee: not saying, ``Depart into the fire prepared for you,'' but ``prepared for the devil.''\endnote{John Chrysostom, \emph{Homilies on the Gospel of Saint Matthew} 23.9, trans.\ G. Prevost, rev.\ M. B. Riddle, \emph{Nicene and Post-Nicene Fathers}, 1st series, vol.~10.}
\end{quote}

Hell, then, was not made for man. It was prepared for the angels who would not serve, and a man comes into it only by joining them.

\subsection{What the Church confesses}

The Church has confessed this sentence in her creeds, in her councils, and in the professions and definitions of her popes. The ancient creed that bears the name of Athanasius ends with the judgment: ``At whose coming all men shall rise again with their bodies, and shall give account for their own works. And they that have done good shall go into life everlasting, and they that have done evil into everlasting fire.''\endnote{The \latin{Quicumque} (Athanasian Creed), DS 76; English from the \emph{Catholic Encyclopedia}, vol.~2 (New York, 1907), s.v.\ ``Athanasian Creed,'' New Advent, accessed 16 October 2026.} The Fourth Lateran Council professed that all, the reprobate and the elect alike, will rise with their own bodies which they now bear, to receive according to their works: \latin{illi cum diabolo poenam perpetuam, et isti cum Christo gloriam sempiternam}, the former perpetual punishment with the devil, the latter everlasting glory with Christ.\endnote{Fourth Lateran Council (1215), constitution 1, \latin{Firmiter}, DS 801.} The Second Council of Lyons professed that the souls of those who die in mortal sin go down at once into hell, \latin{poenis tamen disparibus puniendas}, to be punished, but with unequal punishments;\endnote{Second Council of Lyons (1274), profession of faith of Michael Palaeologus, DS 858; repeated by the Council of Florence (1439), DS 1306.} and Pope Benedict XII defined that the souls of those who die in actual mortal sin go down into hell immediately after death, \latin{ubi poenis infernalibus cruciantur}, where they are tormented with the pains of hell.\endnote{Benedict XII, \latin{Benedictus Deus} (1336), DS 1002. The constitution is expounded in the second volume of this series, \emph{What Death Does Not Finish}.} In 1968, at the close of a Year of Faith, Paul VI pronounced the Church's creed anew: Christ will come to judge the living and the dead, ``those who have responded to the love and piety of God going to eternal life, those who have refused them to the end going to the fire that is not extinguished.''\endnote{Paul VI, \emph{Credo of the People of God} (30 June 1968), n.~12, official Vatican English, vatican.va, accessed 16 October 2026.}

The \emph{Catechism of the Catholic Church} gathers the doctrine in a few sentences. ``To die in mortal sin without repenting and accepting God's merciful love means remaining separated from him for ever by our own free choice. This state of definitive self-exclusion from communion with God and the blessed is called `hell.'\,'' ``Immediately after death the souls of those who die in a state of mortal sin descend into hell, where they suffer the punishments of hell, `eternal fire.' The chief punishment of hell is eternal separation from God, in whom alone man can possess the life and happiness for which he was created and for which he longs.''\endnote{\emph{Catechism of the Catholic Church} 1033 and 1035, official Vatican English, vatican.va, accessed 16 October 2026.}

Together these texts say that hell is real and that those who die in mortal sin enter it at death; that it is eternal; that its chief pain is the loss of God, and that with that loss there is a pain of sense, which Christ and the creeds call fire; that the punishments are unequal, measured to each one's sins; and that at the resurrection the body rises and shares the punishment of the soul. The Church holds each of these with the weight she has given it, and none of them is hers to change.

The \emph{Catechism} gives the reason the Church sets these truths before her children. ``The affirmations of Sacred Scripture and the teachings of the Church on the subject of hell,'' it says, ``are a call to the responsibility incumbent upon man to make use of his freedom in view of his eternal destiny. They are at the same time an urgent call to conversion.''\endnote{\emph{Catechism of the Catholic Church} 1036.} The saints received them in that way. They looked at hell in order not to go there, and they looked long, because a short look is soon forgotten. Their way of looking, which the Church has made her own, was to take each part of the sentence in turn, to let it become real to the mind and heavy to the heart, and then to turn toward the mercy that still stands open.
